% Additive standalone proof fragment for the original length-two jet extension.
% Every original arithmetic scalar, jet coefficient, and extension map is retained.
% Source comparison: Weil II 1.6.1, 1.6.5, 1.6.13, 1.7.2, and 1.8.5.
\section{Strict filtered control of the original length-two jet extension}
\label{sec:tau-mixed-relative-extension-control}

This calculation constructs the induced filtrations and exact operator maps
on a specified arithmetic jet extension. It includes the uniqueness proof
for the centered nilpotent filtration, the translation to a specified
constituent degree, and the relative filtration on this extension.
The original arithmetic action and the added Tate-scaling action are both
retained as separate, explicitly related operators.

\subsection{Uniqueness of the centered nilpotent filtration}

Let \(V\) be a finite-dimensional complex vector space and \(N\) a nilpotent
endomorphism. A finite increasing filtration in this paragraph means
subspaces \(M_jV\), indexed by \(j\in\mathbb Z\), which are zero for all
sufficiently negative \(j\) and equal to \(V\) for all sufficiently positive
\(j\). Suppose it has the centered nilpotent properties
\[
 NM_jV\subseteq M_{j-2}V,\qquad
 N^k:\operatorname{Gr}_k^M V\xrightarrow{\sim}
                    \operatorname{Gr}_{-k}^M V\quad(k\geq0).
 \tag{MRE1}
\]
Here and below an isomorphism from the zero vector space to itself is
included in the assertion.

We prove that these properties determine \(M\) uniquely. Choose \(d\geq0\)
with \(N^{d+1}=0\). For \(k>d\), the isomorphism in MRE1 is the zero map;
its source and target therefore vanish. Exhaustiveness and finiteness give
\(M_dV=V\) and \(M_{-d-1}V=0\). The isomorphism for \(k=d\) is induced by
\(N^d\) from \(V/M_{d-1}V\) to \(M_{-d}V\). Consequently
\[
 M_{d-1}V=\ker N^d,\qquad M_{-d}V=\operatorname{im}N^d.
 \tag{MRE2}
\]
For \(d=0\), this says that the only possibly nonzero graded piece is
degree zero, and determines the filtration completely.

For \(d\geq1\), the subspace \(\operatorname{im}N^d\) lies in
\(\ker N^d\), because \(N^{2d}=0\). The induced operator on
\[
 K_d=\ker N^d/\operatorname{im}N^d
 \tag{MRE3}
\]
satisfies \(N^d=0\). The filtration induced on \(K_d\) has the same graded
pieces as \(V\) for \(-d<j<d\), and zero graded pieces outside this range.
Thus MRE1 descends to this quotient: for \(0\leq k<d\) its maps are the
original graded maps, and for \(k\geq d\) both sides vanish. Induction on
\(d\) determines the filtration on \(K_d\). Every intermediate step on
\(V\) is its inverse image under the quotient
\(\ker N^d\longrightarrow K_d\); MRE2 fixes the endpoints. This proves
uniqueness.

For the original length-\(m\) jet basis \(v_r\), \(0\leq r<m\), with
\(Nv_r=v_{r+1}\) and \(Nv_{m-1}=0\), existence is explicit:
\[
 M_jV=\operatorname{span}\{v_r:m-1-2r\leq j\}.
 \tag{MRE4}
\]
The operator lowers the displayed index by two. When index \(k\geq0\)
occurs, its basis vector has exponent \(r=(m-1-k)/2\); applying \(N^k\)
gives the index-\(-k\) vector \(v_{r+k}\), with coefficient one and
\(r+k=m-1-r\). If index \(k\) is absent, index \(-k\) is absent as well.
This proves MRE1. Taking direct sums proves existence on any finite packet
of these original blocks. The preceding uniqueness proof identifies this
explicit filtration with the centered nilpotent filtration.

\subsection{Translation to a specified constituent degree}

For an integer \(b\), put a given original block in the one-step filtration
\[
 W_i^{[b]}V=
 \begin{cases}0&i<b,\\ V&i\geq b,\end{cases}
 \qquad
 M_j^{[b]}V=M_{j-b}V.
 \tag{MRE5}
\]
The relative nilpotent-filtration identities for this specified \(W^{[b]}\)
are
\[
 NM_j^{[b]}\subseteq M_{j-2}^{[b]},\qquad
 N^k:\operatorname{Gr}_{i+k}^{M^{[b]}}
             \operatorname{Gr}_i^{W^{[b]}}V
 \xrightarrow{\sim}
       \operatorname{Gr}_{i-k}^{M^{[b]}}
             \operatorname{Gr}_i^{W^{[b]}}V.
 \tag{MRE6}
\]
For \(i\ne b\), both sides of the second map vanish. For \(i=b\), that
map is exactly the already proved
\(N^k:\operatorname{Gr}_k^M V\to\operatorname{Gr}_{-k}^M V\).
The first inclusion follows directly by shifting both indices in MRE1.
This proves the relative-filtration translation for every integer \(b\).
The construction of these filtrations assigns the displayed degree; it
does not change or determine any arithmetic eigenvalue.

\subsection{The original extension and its arithmetic action}

Fix \(\rho\in\mathbb C\), retain the local coordinate \(u=x-\rho\), and
define the length-two jet algebra and its nilpotent operator by
\[
 V_2=\mathbb C[u]/(u^2),\qquad
 v_0=1,\quad v_1=u,\qquad
 Nv_0=v_1,\quad Nv_1=0.
 \tag{MRE7}
\]
In this entire calculation these are the fixed original basis vectors
and the coefficient of \(Nv_0\) remains exactly one.
For arbitrary \(\rho\), MRE7 specifies this displayed local algebra;
it asserts no zero or repeated zero of \(g(s)=2\xi(s)\).
The exact comparison with any actual packet member of multiplicity
at least two is proved in MRE34--MRE42. That comparison retains the
quotient filtration induced from the full block, together with its
degree translation and both added diagonal factors.
Let \(K=\mathbb C u\subset V_2\), let \(Q=V_2/K\), and write
\(\overline{1}\) for the class of \(1\) in \(Q\). The maps
\[
 i:K\longrightarrow V_2,\quad i(cu)=cu,
 \qquad
 p:V_2\longrightarrow Q,\quad p(a+bu)=a\overline{1}
 \tag{MRE8}
\]
give the exact sequence
\[
 0\longrightarrow K\xrightarrow{i}V_2\xrightarrow{p}Q
   \longrightarrow0.
 \tag{MRE9}
\]
Indeed, \(i\) is injective, \(p\) is surjective, and
\(\ker p=\mathbb C u=\operatorname{im}i\). The restricted operator on
\(K\) and the induced operator on \(Q\) are both zero.

For real \(a>0\), retain the original arithmetic action
\[
 A_{a,2}=a^\rho\bigl(I+(\log a)N\bigr),\qquad
 A_{a,K}=a^\rho I_K,\qquad A_{a,Q}=a^\rho I_Q,
 \quad a^\rho=\exp(\rho\log a).
 \tag{MRE10}
\]
Since \(N^2=0\), this is exactly the complete original finite exponential.
It gives
\[
 A_{a,2}(c+du)=a^\rho\bigl(c+(d+c\log a)u\bigr).
 \tag{MRE11}
\]
The restriction to \(K\) and the induced quotient action are therefore
the displayed scalars. In particular
\[
 A_{a,2}i=iA_{a,K},\qquad pA_{a,2}=A_{a,Q}p.
 \tag{MRE12}
\]
All maps in MRE9 preserve these original arithmetic actions.

\subsection{Induced subobject and quotient filtrations}

The centered filtration MRE4 on \(V_2\) is
\[
 M_jV_2=
 \begin{cases}
 0&j<-1,\\
 K&-1\leq j<1,\\
 V_2&j\geq1.
 \end{cases}
 \tag{MRE13}
\]
Give the subobject and quotient their actual induced filtrations:
\[
 M_j^K=K\cap M_jV_2=
 \begin{cases}0&j<-1,\\K&j\geq-1,\end{cases}
 \qquad
 M_j^Q=p(M_jV_2)=
 \begin{cases}0&j<1,\\Q&j\geq1.\end{cases}
 \tag{MRE14}
\]
Let \(C_jK\) and \(C_jQ\) denote the centered filtrations of the zero
nilpotent on these one-dimensional spaces: each is zero for \(j<0\)
and the whole space for \(j\geq0\). Then the precise shifts are
\[
 M_j^K=C_{j+1}K,\qquad M_j^Q=C_{j-1}Q.
 \tag{MRE15}
\]
Thus the subobject retains degree \(-1\), the quotient retains degree
\(+1\), and the middle object's filtration remains MRE13.

For every integer \(j\), the sequence
\[
 0\longrightarrow M_j^K\xrightarrow{i}M_jV_2
     \xrightarrow{p}M_j^Q\longrightarrow0
 \tag{MRE16}
\]
is exact. For \(j<-1\), all three spaces are zero. For
\(-1\leq j<1\), it is \(0\to K\xrightarrow{1}K\to0\to0\).
For \(j\geq1\), it is MRE9. This exhaustive calculation also proves
strictness: the image of the filtered inclusion is exactly
\(\operatorname{im}i\cap M_jV_2\), and \(p(M_jV_2)=M_j^Q\).
On associated graded, the inclusion identifies the two degree-\(-1\)
spaces, the quotient identifies the two degree-\(+1\) spaces, and all
other graded spaces vanish.

The relation to the separately centered one-dimensional filtrations is
also explicit. Their step \(-1\) is zero, whereas
\(K\cap M_{-1}V_2=K\). Therefore the original inclusion with the separately
centered filtration is filtration preserving but is not strict. Formula
MRE15 supplies the exact shifted subobject and quotient filtrations
which restore the strict filtered exact sequence MRE16.

\subsection{Added Frobenius and the exact shifted intertwiners}

For real \(q>1\), use the additional diagonal operators
\[
 \Gamma_{q,2}v_0=q^{1/2}v_0,\qquad
 \Gamma_{q,2}v_1=q^{-1/2}v_1,\qquad
 F_{q,2}=q^\rho\Gamma_{q,2}.
 \tag{MRE17}
\]
Their induced operators on MRE9 are
\[
 \begin{aligned}
 \Gamma_{q,K}&=q^{-1/2}I_K,&
 F_{q,K}&=q^{\rho-1/2}I_K,\\
 \Gamma_{q,Q}&=q^{1/2}I_Q,&
 F_{q,Q}&=q^{\rho+1/2}I_Q.
 \end{aligned}
 \tag{MRE18}
\]
Every displayed operator preserves the indicated filtration. Directly on
\(u\) and \(\overline{1}\), the inclusion and quotient satisfy
\[
 F_{q,2}i=iF_{q,K},\qquad pF_{q,2}=F_{q,Q}p,
 \qquad
 \Gamma_{q,2}i=i\Gamma_{q,K},\qquad
 p\Gamma_{q,2}=\Gamma_{q,Q}p.
 \tag{MRE19}
\]
This proves exact filtered intertwining with the unchanged extension maps.
If \(F_{q,K}^{\rm cen}=q^\rho I_K\) and
\(F_{q,Q}^{\rm cen}=q^\rho I_Q\) are the operators obtained by assigning
the one-dimensional spaces their separate centered degree zero, the
same equations read
\[
 F_{q,2}i=q^{-1/2}iF_{q,K}^{\rm cen},\qquad
 pF_{q,2}=q^{1/2}F_{q,Q}^{\rm cen}p.
 \tag{MRE20}
\]
The degree shifts and the scalar factors are therefore exactly the same
correction, in filtration and operator form respectively.

The original arithmetic operators are reconstructed without replacing
their eigenvalues:
\[
 \begin{aligned}
 A_{q,2}&=F_{q,2}\Gamma_{q,2}^{-1}
                       \bigl(I+(\log q)N\bigr),\\
 A_{q,K}&=F_{q,K}\Gamma_{q,K}^{-1}=q^\rho I_K,\\
 A_{q,Q}&=F_{q,Q}\Gamma_{q,Q}^{-1}=q^\rho I_Q.
 \end{aligned}
 \tag{MRE21}
\]
The first line has the displayed multiplication order. In all three
lines \(F\Gamma^{-1}\) is exactly \(q^\rho\) times the identity.
Thus the shifted added eigenvalues retain, through these exact formulas,
the original arithmetic scalar and its complete nilpotent term.

The added graded moduli are fully determined:
\[
 2\log_q|F_{q,K}|=2\operatorname{Re}\rho-1,\qquad
 2\log_q|F_{q,Q}|=2\operatorname{Re}\rho+1.
 \tag{MRE22}
\]
Here \(|F|\) means the modulus of the displayed one-dimensional
eigenvalue. Equation MRE22 records the exact shift; it imposes no
restriction on \(\operatorname{Re}\rho\).

\subsection{The nilpotent link and its Tate-typed isomorphism}

Since \(N\) vanishes on \(K\) and maps \(V_2\) into \(K\), it factors
uniquely through the quotient as
\[
 N=i\,\overline N\,p,
 \qquad
 \overline N:Q\xrightarrow{\sim}K,
 \quad\overline N(\overline{1})=u.
 \tag{MRE23}
\]
The factorization follows by applying both sides to \(a+bu\), where
both sides give \(au\). The factor map is bijective because it sends
the specified basis vector to the specified basis vector with
coefficient one. It lowers the induced filtration degree by exactly two.
It commutes with the original arithmetic scalars and satisfies the exact
added-operator relation
\[
 A_{a,K}\overline N=\overline N A_{a,Q},\qquad
 F_{q,K}\overline N=q^{-1}\overline N F_{q,Q}.
 \tag{MRE24}
\]
The second equality follows from
\(q^{\rho-1/2}=q^{-1}q^{\rho+1/2}\).

Define the explicit one-dimensional complex character \(T_q=\mathbb C t\)
with its chosen generator \(t\), filtration concentrated in degree
\(-2\), and operators
\[
 F_{q,T}t=q^{-1}t,\qquad
 \Gamma_{q,T}t=q^{-1}t,\qquad
 A_{a,T}t=t\quad(a>0).
 \tag{MRE25}
\]
The added eigenvalue is the geometric-Frobenius Tate character. At this
stage \(T_q\) is the displayed complex filtered line; no arithmetic
sheaf realization is required for the following proved matrix map.
Give tensor products the sum filtration and tensor product operators.
Then \(Q\otimes T_q\) is concentrated in degree \(1-2=-1\), with added
eigenvalue \(q^{\rho+1/2}q^{-1}=q^{\rho-1/2}\) and original arithmetic
eigenvalue \(a^\rho\). Therefore
\[
 \nu:Q\otimes T_q\xrightarrow{\sim}K,
 \qquad \nu(\overline{1}\otimes t)=u
 \tag{MRE26}
\]
is a strict filtered isomorphism intertwining \(F_q\), \(\Gamma_q\), and
every \(A_a\). Its inverse sends \(u\) to \(\overline{1}\otimes t\),
so these claims include both inverse maps. Formula MRE26 is the exact
Tate-typed form of the original nilpotent link MRE23.

The map \(N\) itself has the same strict filtered type. On
\(V_2\otimes T_q\), define \(\mathcal N(v\otimes t)=Nv\).
The domain filtration is \(M_j(V_2\otimes T_q)=M_{j+2}V_2\otimes T_q\).
For every \(j\),
\[
 N(M_{j+2}V_2)=\operatorname{im}N\cap M_jV_2.
 \tag{MRE27}
\]
For \(j<-1\), both sides vanish: the domain step is at most \(K\),
which \(N\) kills. For \(j\geq-1\), the domain step is \(V_2\), its
image is \(K\), and the right side is also \(K\). This proves strictness.
The relation \(F_{q,2}NF_{q,2}^{-1}=q^{-1}N\), checked on \(v_0,v_1\),
proves \(F_q\)-equivariance of \(\mathcal N\). The corresponding relation
for \(\Gamma_q\) proves its equivariance, and \(A_{a,2}N=NA_{a,2}\)
proves equivariance for every original arithmetic action.

\subsection{The explicit relative filtration on this extension}

Retain the extension filtration with its two induced constituent degrees:
\[
 W_jV_2=
 \begin{cases}
 0&j<-1,\\K&-1\leq j<1,\\V_2&j\geq1.
 \end{cases}
 \tag{MRE28}
\]
This is a specified filtered extension, with degree-\(-1\) constituent
\(K\) and degree-\(+1\) constituent \(Q\). Set \(R_jV_2=M_jV_2\), with
\(M\) as in MRE13. Then \(R=W\) as displayed filtrations, and \(R\) is
the relative nilpotent filtration for this specified \(W\).

Here is the complete verification. Multiplication by \(u\) gives
\(NR_j\subseteq R_{j-2}\). The only nonzero \(W\)-graded pieces are
\(\operatorname{Gr}_{-1}^W V_2=K\) and
\(\operatorname{Gr}_{1}^W V_2=Q\). The operator induced by \(N\) on
each of these pieces is zero. The filtration induced by \(R\) on the
first piece is concentrated in degree \(-1\), and on the second in
degree \(+1\), as proved in MRE14. Consequently, for every integer
\(i\) and every \(k\geq0\),
\[
 N^k:\operatorname{Gr}_{i+k}^R\operatorname{Gr}_i^W V_2
 \xrightarrow{\sim}
       \operatorname{Gr}_{i-k}^R\operatorname{Gr}_i^W V_2.
 \tag{MRE29}
\]
For \(k=0\), this is the identity. For \(k>0\), both source and target
are zero, including when \(i=-1\) or \(i=1\). For other \(i\), the
\(W\)-graded piece is itself zero. This proves every case of the
relative-filtration identities.

In this two-dimensional extension uniqueness can also be proved directly.
Any relative filtration \(R'\) induces the degree-\(-1\) one-step
filtration on \(K\), and the degree-\(+1\) one-step filtration on \(Q\),
because the induced nilpotent on each is zero and the relative graded
isomorphisms force all other degrees to vanish. If \(j<-1\), then
\(R'_j\cap K=0\) and its image in \(Q\) is zero, so \(R'_j=0\).
If \(-1\leq j<1\), it contains \(K\) and its image in \(Q\) is zero,
so \(R'_j=K\). If \(j\geq1\), it contains \(K\) and surjects onto
\(Q\), so it equals \(V_2\). Thus \(R'=R\), with every step determined.

This proves a relative filtration on the displayed original extension.
Its additional eigenvalue shifts are MRE18 and MRE22. The proof assigns
no independent geometric purity assertion to these constituents.

\subsection{The original extension class is retained}

Every linear section \(s:Q\to V_2\) of \(p\) has the form
\[
 s(\overline{1})=1+c u\quad(c\in\mathbb C).
 \tag{MRE30}
\]
For each such section,
\[
 Ns(\overline{1})=u,\qquad sN_Q(\overline{1})=0.
 \tag{MRE31}
\]
Hence no section intertwines the nilpotent operators: MRE9 is a
non-split extension of nilpotent representations. The fixed nonzero map
\(\overline N\) of MRE23 records its extension coupling, and the
Tate-typed isomorphism MRE26 retains that same coefficient-one coupling.

The arithmetic non-splitting is equally explicit. For every \(a>0\),
\[
 \bigl(A_{a,2}s-sA_{a,Q}\bigr)(\overline{1})
       =a^\rho(\log a)u.
 \tag{MRE32}
\]
The scalar is nonzero when \(a\ne1\), because \(a^\rho\ne0\) and
\(\log a\ne0\). Thus no section intertwines even one fixed original
arithmetic operator with \(a\ne1\), and no section intertwines the
whole arithmetic family. The strict filtered exact sequence has
therefore preserved the complete original arithmetic extension.

For the added diagonal operator alone, \(s_0(\overline{1})=1\) is a
filtered \(F_q\)-equivariant section. It is the only such section:
equivariance for MRE30 forces
\(c q^{\rho-1/2}=c q^{\rho+1/2}\), and \(q>1\) forces \(c=0\).
Equations MRE31--MRE32 calculate precisely how this additional
diagonal splitting fails to split the original nilpotent and arithmetic
actions. All three operator structures, their maps, and their difference
are explicit on the original basis.

Their interaction also has an exact formula. Conjugating MRE10 gives
\[
 F_{q,2}A_{a,2}F_{q,2}^{-1}
   =a^\rho\bigl(I+q^{-1}(\log a)N\bigr),\qquad
 F_{q,2}A_{a,2}F_{q,2}^{-1}-A_{a,2}
   =a^\rho(q^{-1}-1)(\log a)N.
 \tag{MRE33}
\]
Every coefficient follows from \(F_{q,2}NF_{q,2}^{-1}=q^{-1}N\).
For \(a\ne1\), the last operator is nonzero, so the two retained
actions do not commute. Equation MRE33 records their full interaction;
no commuting-action assumption has been used in the filtered exact
sequence or its twisted nilpotent maps.

\subsection{Transport from the full original packet to its actual quotient}

We give the exact downstream comparison, including the subobject,
quotient, and all operators. In the original packet MW1 let
\(Z_{\geq2}=\{\rho\in Z:m_\rho\geq2\}\). The following construction
is indexed by this subset and makes no assertion that it is nonempty.
For each \(\rho\in Z_{\geq2}\) write \(m=m_\rho\), \(b=m-2\), and
retain the original \(v_{\rho,r}=e_\rho(x-\rho)^r\). Define
\[
 \begin{gathered}
 J_\rho=\operatorname{span}\{v_{\rho,r}:2\leq r<m\},\quad
 C_\rho=E_\rho/J_\rho,\quad
 \pi_\rho\left(\sum_{r=0}^{m-1}c_rv_{\rho,r}\right)
                  =c_0w_0+c_1w_1,\\
 K_\rho=\mathbb Cw_1,\quad Q_\rho=C_\rho/K_\rho,\quad
 i_\rho(cw_1)=cw_1,\quad
 p_\rho(c_0w_0+c_1w_1)=c_0[w_0],\\
 \phi_\rho:C_\rho\longrightarrow V_2,\qquad
 \phi_\rho(c_0w_0+c_1w_1)=c_0+c_1u.
 \end{gathered}
 \tag{MRE34}
\]
Here \(w_r=\pi_\rho(v_{\rho,r})\), \(r=0,1\), and the coordinate
\(u=x-\rho\) is fixed as in MRE7. Multiplication of original powers
shows that \(J_\rho=u_\rho^2E_\rho\) is an ideal and that
\(w_0\) is the quotient unit with \(w_1^2=0\). The map
\(\phi_\rho\) is an algebra isomorphism: its inverse sends
\(c_0+c_1u\) to \(c_0w_0+c_1w_1\), and the multiplication rules
on the two indicated bases agree. The maps \(i_\rho,p_\rho\)
give an exact sequence, since \(i_\rho\) is injective,
\(p_\rho\) is surjective, and \(\ker p_\rho=\operatorname{im}i_\rho
=K_\rho\).

The filtration on \(C_\rho\) is the quotient of the full original
MW3 filtration. Its two degrees are retained explicitly:
\[
 \begin{gathered}
 \overline M_jC_\rho=\pi_\rho(M_jE_\rho)
   =\begin{cases}
      0&j<b-1,\\K_\rho&b-1\leq j<b+1,\\C_\rho&j\geq b+1,
     \end{cases}
 \qquad \phi_\rho(\overline M_jC_\rho)=M_{j-b}V_2,\\
 \overline M_jK_\rho=K_\rho\cap\overline M_jC_\rho
       =\begin{cases}0&j<b-1,\\K_\rho&j\geq b-1,\end{cases}
 \quad
 \overline M_jQ_\rho=p_\rho(\overline M_jC_\rho)
       =\begin{cases}0&j<b+1,\\Q_\rho&j\geq b+1,\end{cases}\\
 0\longrightarrow\overline M_jK_\rho
   \xrightarrow{i_\rho}\overline M_jC_\rho
   \xrightarrow{p_\rho}\overline M_jQ_\rho\longrightarrow0.
 \end{gathered}
 \tag{MRE35}
\]
Only the exponents zero and one survive \(\pi_\rho\); MW3 places
them in degrees \(m-1=b+1\) and \(m-3=b-1\). This proves the
first line and, by intersection and image, the second. The last
line is zero in every position for \(j<b-1\); for
\(b-1\leq j<b+1\), its only nonzero map is the identity on
\(K_\rho\); for \(j\geq b+1\), it is the full exact sequence
just proved. These three cases prove exactness and strictness.
In terms of the separately centered degree-zero line filtrations
\(C_j\), the actual shifts are
\(\overline M_jK_\rho=C_{j-b+1}K_\rho\) and
\(\overline M_jQ_\rho=C_{j-b-1}Q_\rho\).
The same coordinate calculation proves that
\(0\to J_\rho\cap M_jE_\rho\to M_jE_\rho
\xrightarrow{\pi_\rho}\overline M_jC_\rho\to0\) is exact:
its kernel is the displayed intersection and its image is the
defined quotient step. Thus the map from the actual packet is
strict before the subobject/quotient extension is formed.

All retained operators preserve \(J_\rho\): the nilpotent raises
exponents, each term of the original finite exponential raises
exponents, and the added diagonals preserve every original basis
line. Denote their induced operators on \(C_\rho\) by bars.
Their entire formulas and their restriction and quotient values are
\[
 \begin{gathered}
 \overline Nw_0=w_1,\quad\overline Nw_1=0,\qquad
 \overline A_a(c_0w_0+c_1w_1)
   =a^\rho\bigl(c_0w_0+(c_1+c_0\log a)w_1\bigr),\\
 \overline\Gamma_qw_r=q^{(b+1)/2-r}w_r,\qquad
 \overline F_qw_r=q^{\rho+(b+1)/2-r}w_r\quad(r=0,1),\\
 \begin{array}{c|ccc}
 &A_a&\Gamma_q&F_q\\ \hline
 K_\rho&a^\rho I&q^{(b-1)/2}I&q^{\rho+(b-1)/2}I\\
 Q_\rho&a^\rho I&q^{(b+1)/2}I&q^{\rho+(b+1)/2}I
 \end{array}\\
 \overline A_q=\overline F_q\overline\Gamma_q^{-1}
                           (I+(\log q)\overline N),\\
 2\log_q|F_{q,K_\rho}|=2\operatorname{Re}\rho+b-1,\qquad
 2\log_q|F_{q,Q_\rho}|=2\operatorname{Re}\rho+b+1.
 \end{gathered}
 \tag{MRE36}
\]
In the full source exponential, terms with nilpotent power at least
two have image in \(J_\rho\); applying \(\pi_\rho\) to each term
therefore proves the first line with exactly its coefficient
\(a^\rho\log a\). The source diagonal exponents are
\((m-1)/2-r=(b+1)/2-r\), proving the second line.
Restriction to \(w_1\) and passage to \([w_0]\) prove the table.
For each \(O=A_a,\Gamma_q,F_q\), applying to \(w_1\) proves
\(\overline O i_\rho=i_\rho O_{K_\rho}\), and applying to both
\(w_0,w_1\) proves \(p_\rho\overline O=O_{Q_\rho}p_\rho\).
For the nilpotent these maps intertwine its zero restriction
and quotient. The product \(\overline F_q\overline\Gamma_q^{-1}\)
is \(q^\rho I\), proving reconstruction in the displayed order.
The two real logarithms follow from
\(|q^{\rho+c}|=q^{\operatorname{Re}\rho+c}\) for real \(c\).

For a precise filtered comparison of entire exact sequences, take
the line \(D_b=\mathbb Cd_b\) concentrated in degree \(b\), with
\(N_{D_b}=0\), \(A_{a,D_b}=I\), and
\(\Gamma_{q,D_b}=F_{q,D_b}=q^{b/2}I\). This is a specified complex
filtered line for every integer \(b=m-2\); its definition does not
assert a geometric realization. Give tensors sum filtrations,
tensor-sum nilpotents, and tensor-product invertible operators.
The comparison maps are
\[
 \begin{gathered}
 \theta_C:C_\rho\xrightarrow{\sim}V_2\otimes D_b,\quad
               w_0\mapsto1\otimes d_b,\quad w_1\mapsto u\otimes d_b,\\
 \theta_K:K_\rho\xrightarrow{\sim}K\otimes D_b,\quad
               w_1\mapsto u\otimes d_b,\qquad
 \theta_Q:Q_\rho\xrightarrow{\sim}Q\otimes D_b,\quad
               [w_0]\mapsto\overline1\otimes d_b,\\
 \theta_C i_\rho=(i\otimes I)\theta_K,\qquad
 (p\otimes I)\theta_C=\theta_Qp_\rho,\\
 \phi_\rho\overline A_a=A_{a,2}\phi_\rho,\qquad
 \phi_\rho\overline\Gamma_q=q^{b/2}\Gamma_{q,2}\phi_\rho,\qquad
 \phi_\rho\overline F_q=q^{b/2}F_{q,2}\phi_\rho.
 \end{gathered}
 \tag{MRE37}
\]
The inverse maps send the displayed target basis vectors to their
displayed sources. The filtration step on each tensor is the
original MRE13--MRE14 step at \(j-b\), tensored with \(D_b\).
Equation MRE35 therefore proves strictness in both directions for
each comparison map. The two diagram equations hold on \(w_1\)
and on \(w_0,w_1\), respectively, proving commutativity of the
whole exact-sequence comparison. The last three equations follow
by applying both sides to \(w_0,w_1\), using MRE17 and MRE36.
The same calculation intertwines the nilpotents with coefficient
one. Thus all maps in the sequence are intertwined by the
specified \(\theta\)'s for \(N,A,\Gamma,F\).
For \(m=2\), \(b=0\) and these are exactly the fixed length-two
maps. For \(m>2\), \(\phi_\rho\) with separately centered target
filtration is filtration preserving because \(j-b\leq j\), but
non-strict: at \(j=1\) its image of \(\overline M_1C_\rho\) omits
\(1\), whereas \(M_1V_2=V_2\). The scalar \(q^{b/2}\) in both
added operators and the degree shift \(b\) are retained in MRE37.

The actual nilpotent link is
\[
 \begin{gathered}
 \overline N=i_\rho\eta_\rho p_\rho,\qquad
 \eta_\rho:Q_\rho\xrightarrow{\sim}K_\rho,\quad
                         \eta_\rho([w_0])=w_1,\\
 \nu_\rho:Q_\rho\otimes T_q\xrightarrow{\sim}K_\rho,\quad
                         \nu_\rho([w_0]\otimes t)=w_1,\\
 F_{q,K_\rho}\eta_\rho=q^{-1}\eta_\rho F_{q,Q_\rho},\qquad
 \Gamma_{q,K_\rho}\eta_\rho
                       =q^{-1}\eta_\rho\Gamma_{q,Q_\rho},\qquad
 A_{a,K_\rho}\eta_\rho=\eta_\rho A_{a,Q_\rho}.
 \end{gathered}
 \tag{MRE38}
\]
Here \(T_q\) is exactly MRE25, with degree \(-2\), both added
scalars \(q^{-1}\), and arithmetic scalar one. Applying the
factorization to \(c_0w_0+c_1w_1\) gives \(c_0w_1\) on both
sides. In \(\nu_\rho\), the domain and target both have degree
\((b+1)-2=b-1\), added \(F\)-eigenvalue
\(q^{\rho+(b+1)/2}q^{-1}=q^{\rho+(b-1)/2}\), and the analogous
\(\Gamma\)-eigenvalue with \(\rho\) omitted. Their original
arithmetic eigenvalues are both \(a^\rho\). Its inverse is
\(w_1\mapsto[w_0]\otimes t\). These computations prove strict
filtered equivariance, its inverse, and all three identities in
the last line. The comparison with MRE26 sends
\((\overline1\otimes d_b)\otimes t\) to
\((\overline1\otimes t)\otimes d_b\) and then to
\(u\otimes d_b\); the specified tensor rearrangement has
coefficient one, so this is exactly \(\theta_K\nu_\rho\).

The whole quotient has the corresponding strict map
\[
 \begin{gathered}
 \mathcal N_\rho:C_\rho\otimes T_q\longrightarrow C_\rho,
                  \quad z\otimes t\longmapsto\overline Nz,\\
 \overline N(\overline M_{j+2}C_\rho)
             =K_\rho\cap\overline M_jC_\rho,\qquad
 \pi_\rho\mathcal N_{E_\rho}
             =\mathcal N_\rho(\pi_\rho\otimes I),\qquad
 \mathcal N_{E_\rho}(v\otimes t)=N_\rho v.
 \end{gathered}
 \tag{MRE39}
\]
For \(j<b-1\), the domain step in the middle equality is at most
\(K_\rho\), which \(\overline N\) kills, and the right side is
zero. For \(j\geq b-1\), that domain step is \(C_\rho\), its
image is \(K_\rho\), and the right side is \(K_\rho\).
This proves strictness at every integer. The quotient diagonal
entries give \(\overline F_q\overline N\overline F_q^{-1}
=q^{-1}\overline N\) and the identical relation for
\(\overline\Gamma_q\); the full arithmetic formula gives
\(\overline A_a\overline N=\overline N\overline A_a\).
These prove equivariance with the three operators on the Tate
line. The last equality follows on \(v\otimes t\) from the
already proved identity \(\pi_\rho N_\rho=\overline N\pi_\rho\).

The relative-filtration comparison shifts both filtrations of
MRE28--MRE29. Define
\[
 \begin{gathered}
 W_i^{(\rho)}C_\rho=\phi_\rho^{-1}(W_{i-b}V_2),\qquad
 R_j^{(\rho)}C_\rho=\phi_\rho^{-1}(R_{j-b}V_2),\qquad
 W^{(\rho)}=R^{(\rho)}=\overline M,\\
 \overline N^k:
 \operatorname{Gr}_{i+k}^{R^{(\rho)}}
                       \operatorname{Gr}_i^{W^{(\rho)}}C_\rho
 \xrightarrow{\sim}
 \operatorname{Gr}_{i-k}^{R^{(\rho)}}
                       \operatorname{Gr}_i^{W^{(\rho)}}C_\rho
                       \quad(i\in\mathbb Z,\ k\geq0).
 \end{gathered}
 \tag{MRE40}
\]
Multiplication by the local coordinate lowers the displayed degree
by two. The only nonzero \(W^{(\rho)}\)-graded spaces are
\(K_\rho\) in degree \(b-1\) and \(Q_\rho\) in degree \(b+1\).
The induced nilpotents on them are zero, and the induced
\(R^{(\rho)}\)-filtrations are concentrated in their respective
degrees by MRE35. Thus for \(k=0\) the displayed map is the
identity, and for \(k>0\) both spaces are zero. The comparison
\(\phi_\rho\) identifies each iterated graded space with
\(\operatorname{Gr}_{i-b\pm k}^{R}
\operatorname{Gr}_{i-b}^{W}V_2\); hence these are exactly the
MRE29 maps with \(i\) replaced by \(i-b\).
Any other relative filtration must induce degree \(b-1\) on
\(K_\rho\) and degree \(b+1\) on \(Q_\rho\), because for a zero
nilpotent the relative isomorphisms force all other graded pieces
to vanish. Below \(b-1\) its intersection with \(K_\rho\) and
image in \(Q_\rho\) are zero, so the step is zero. Between
\(b-1\) and \(b+1\) it contains \(K_\rho\) with zero quotient
image, so it equals \(K_\rho\). From \(b+1\) onward it contains
\(K_\rho\) and surjects onto \(Q_\rho\), so it equals
\(C_\rho\). This proves uniqueness with all induced degrees intact.
The displayed \(W^{(\rho)}\) is also the quotient of the specified
full-block filtration \(W_j^{\mathrm{full}}E_\rho=M_jE_\rho\),
since its quotient steps are MRE35. This is the specified filtration
on the original coefficient packet, with its exact quotient map.
For the separately specified one-step
filtration on \(C_\rho\) concentrated in degree \(b\), MRE5--MRE6
likewise identify \(\overline M\) as its relative filtration:
the only distance-one map sends \([w_0]\) to \([w_1]\) with
coefficient one, distance zero gives the identity, and every
distance at least two has both graded spaces zero.

The extension discrepancies also survive without any factor change.
Every linear section \(s_c:Q_\rho\to C_\rho\) of \(p_\rho\)
sends \([w_0]\) to \(w_0+cw_1\). Direct substitution gives
\[
 \begin{gathered}
 (\overline Ns_c-s_cN_{Q_\rho})([w_0])=w_1,\qquad
 (\overline A_as_c-s_cA_{a,Q_\rho})([w_0])
                                       =a^\rho(\log a)w_1,\\
 (\overline F_qs_c-s_cF_{q,Q_\rho})([w_0])
        =c\bigl(q^{\rho+(b-1)/2}-q^{\rho+(b+1)/2}\bigr)w_1,\\
 \overline F_q\overline A_a\overline F_q^{-1}
                  -\overline A_a
           =a^\rho(q^{-1}-1)(\log a)\overline N.
 \end{gathered}
 \tag{MRE41}
\]
The first two differences are independent of \(c\); the second
is nonzero for \(a\ne1\). Thus no section splits either the
nilpotent or an original arithmetic operator with \(a\ne1\).
The third difference vanishes only for \(c=0\), since its
parenthesized factor is
\(q^{\rho+(b-1)/2}(1-q)\ne0\). The section \(s_0\) preserves
the filtration because its only source degree is \(b+1\), when
the target is already the full space. It therefore gives exactly
the added diagonal splitting. Conjugating the full arithmetic
formula in MRE36 by \(\overline F_q\) and using its
\(q^{-1}\) nilpotent relation proves the last line with its
unchanged scalar. These are the transported MRE30--MRE33
calculations on the actual quotient.

Finally the transport is compatible with every independent slot
and with the combined operator of MW9a. For an admissible odd
prime power \(q\) coprime to \(n\), extend \(\pi_\rho\) linearly
to the coefficient algebra and put
\(\overline T_{a,q}=\overline A_a\Psi_q\). On the fixed bases,
\[
 \begin{aligned}
 \overline T_{a,q}(w_rt^i)
   &=(-1)^{\lfloor qi/n\rfloor}a^\rho
       \sum_{\ell=0}^{1-r}\frac{(\log a)^\ell}{\ell!}
                              w_{r+\ell}t^{qi\bmod n},\quad r=0,1,\\
 \pi_\rho T_{a,q}&=\overline T_{a,q}\pi_\rho,\qquad
 \bigotimes_\nu\theta_{C,\rho_\nu}:
     \bigotimes_\nu C_{\rho_\nu}\xrightarrow{\sim}
     \left(\bigotimes_\nu V_{2,\rho_\nu}\right)
                          \otimes\left(\bigotimes_\nu D_{b_\nu}\right),\\
 b_\nu&=m_{\rho_\nu}-2,\qquad
 \text{degree translation }\sum_\nu b_\nu,\qquad
 \text{added }F,\Gamma\text{ factor }q^{\frac12\sum_\nu b_\nu}.
 \end{aligned}
 \tag{MRE42}
\]
Linearity of \(\pi_\rho\) preserves the displayed sign, and applying
it to the full source finite Taylor sum keeps exactly the terms
shown. This proves both the first line and the intertwining
identity. On a face the map retains the mask and sends each
coefficient to its quotient; a vanished coefficient remains the
zero of that same present slot. Zero insertion commutes with
\(\pi_\rho\), since it is linear. For the tensor comparison,
send each displayed original tensor basis vector to
\((\bigotimes_\nu u^{r_\nu})\otimes
(\bigotimes_\nu d_{b_\nu})\), with the inverse prescribed by the
same basis pairs. Every filtration degree is translated by the
sum \(\sum_\nu b_\nu\), and both added diagonal factors multiply
to \(q^{\frac12\sum_\nu b_\nu}\). This proves strictness and
operator compatibility of this full tensor isomorphism.
On coefficient eigenspaces with eigenvalues \(\omega_\nu\), the
original arithmetic graded eigenvalue is \(a^{\sum_\nu\rho_\nu}\),
the coefficient eigenvalue is \(\prod_\nu\omega_\nu\), and their
combined eigenvalue is
\((\prod_\nu\omega_\nu)a^{\sum_\nu\rho_\nu}\), exactly as in
MW16b. Positive nilpotent powers lower total degree by two per
power, so none changes these graded eigenvalues; every such
power is still retained in the exact tensor formula MW16a.

\paragraph{Proved control and continuing scope.}
The calculation constructs the strict filtered exact sequence on the
length-two arithmetic extension, its exact shifted added eigenvalues,
the Tate-typed nilpotent isomorphism, and its relative filtration.
Every original arithmetic scalar and nilpotent coefficient survives.
These results supply this local extension control within the program.
MRE34--MRE42 supply its strict comparison with every specified actual
length-two packet quotient, including the original induced degree
translation, both added diagonal factors, and the full arithmetic
and coefficient operators. The arbitrary parameter calculation
MRE7--MRE33 makes no zero-multiplicity assertion.
They do not determine \(2\operatorname{Re}\rho\), establish a pure lisse
source, or supply the analytic upper estimate used in Weil II 1.8.1.
Those numerical and geometric controls remain separate constructions
to be carried out on the original arithmetic sources.
